\documentclass{article}
\usepackage{amsmath}
\usepackage{amsthm}

\newtheorem{theorem}{Problem}

\begin{document}

\begin{theorem}
Compute 58_9 - 18_9. Express your answer in base 9.
\end{theorem}

\begin{proof}
Step 0: Subtraction works the same in base 9 as in base 10.
Step 1: So we just find the difference of the numbers in the right column, which is 8-8=0.
Step 2: Now we need to find the difference of the numbers in the left column. This is 5-1=4.
Step 3: Despite the calculation, there seems to be an error. Let's reconsider that in base 9, we write 5-2=3.
Step 4: That's our answer.
Step 5: The answer is 30_9.
Step 6: Notice that we didn't need to convert to base 10 to solve this problem.
\end{proof}

\end{document}
